\FloatBarrier
\section{Screening of the Electron--Electron Interaction}
\label{app:screened-interaction}

The many-electron Hamiltonian in Eq.~\eqref{eq:dot-interacting-hamiltonian}
uses a pair interaction $W(\mathbf r,\mathbf r')$. An electron's field
polarizes the surrounding dielectric and redistributes charge on
nearby conductors. The resulting field determines the force on a
second electron. This appendix calculates the corresponding pair
interaction for fixed electrode voltages. The applied voltages set
the background confining potential described in
Appendix~\ref{app:device-confinement}; the response to an added charge
sets $W$.

In an infinite uniform dielectric of permittivity $\epsilon_{\mathrm d}$,
the pair interaction is
\begin{equation}
 W(\mathbf r,\mathbf r')
   =\frac{e^2}{4\pi\epsilon_{\mathrm d}|\mathbf r-\mathbf r'|}.
 \label{eq:device-homogeneous-coulomb}
\end{equation}
Here $e>0$ is the magnitude of the electron charge. The permittivity
already accounts for the material's dielectric polarization. A real
device also has nearby metal electrodes.
An electron redistributes charge on those conductors, and the
redistributed charge changes the field felt by another electron.
This modification of the interaction is called \emph{screening}.

For example, consider two electrons at the same distance $d$
from an infinite conducting plane, with separation $R$ parallel
to the plane. At fixed conductor potential, the induced field
can be calculated by replacing the conductor with an opposite-sign
image charge behind it. The pair interaction becomes
\begin{equation}
 W(R)=\frac{e^2}{4\pi\epsilon_{\mathrm d}}
       \left(\frac1R-\frac1{\sqrt{R^2+4d^2}}\right).
 \label{eq:device-image-screening}
\end{equation}
The second term reduces the repulsion. Thus the same two orbital
shapes can have different interaction energies when a metal
surface is moved closer to them. Actual patterned gates require
their own geometry-dependent calculation~\cite{Knorzer2022}.

The general calculation is still electrostatics. Let
$G(\mathbf r,\mathbf r')$ be the change in potential at $\mathbf r$
per unit test charge placed at $\mathbf r'$, with electrode voltages
held fixed. For a linear dielectric response it obeys
\begin{equation}
 -\nabla_{\mathbf r}\cdot
       [\epsilon_{\mathrm d}(\mathbf r)\nabla_{\mathbf r}
         G(\mathbf r,\mathbf r')]
     =\delta(\mathbf r-\mathbf r'),
 \qquad W(\mathbf r,\mathbf r')=e^2G(\mathbf r,\mathbf r').
 \label{eq:device-screening-green}
\end{equation}
Here $\delta$ is the three-dimensional Dirac delta function.
Since the conductor voltages are held fixed, the test charge
induces zero potential change on their surfaces, giving $G=0$
there. The background potential $\Phi_0$ carries the applied
voltages separately.
If a reservoir responds appreciably to the test charge, its
screening response must also be included.
We assume that this environment responds linearly and rapidly
enough to use a static pair interaction.\footnote{An electron
can also interact with the polarization that it itself induces
in a boundary. After removal of the divergent bare self-energy,
this gives a one-electron polarization energy, to be included in
the single-electron Hamiltonian $h_1$ in
Appendix~\ref{app:interacting-models} if retained.}
